\section{关键源代码节选}\label{app:source-code}
以下片段对应本文的四个模型步骤：计算图预处理、候选方案构造、轻量排序和评价接口，仅列出能够对应公式、约束和事件定义的关键实现。

\begingroup
\IfFontExistsTF{Menlo}{\setmonofont{Menlo}}{\setmonofont{Consolas}}
\fvset{fontsize=\scriptsize,baselinestretch=1,breaklines=true,breakanywhere=true,breakanywheresymbolpre={},breaksymbolleft={},numbers=left,numbersep=4pt,xleftmargin=2em}

\subsection{计算图与任务优先级}
\noindent\textbf{计算图与任务优先级：}张量依赖收缩、关键路径和释放优先级。\par\nobreak
\VerbatimInput[firstline=22,lastline=70,firstnumber=22]{support/experiment/v7_fast_20260927/src/huawei_code/graph.py}
\VerbatimInput[firstline=91,lastline=107,firstnumber=91]{support/experiment/v7_fast_20260927/src/huawei_code/graph.py}
\VerbatimInput[firstline=110,lastline=141,firstnumber=110]{support/experiment/v7_fast_20260927/src/huawei_code/graph.py}
\VerbatimInput[firstline=144,lastline=170,firstnumber=144]{support/experiment/v7_fast_20260927/src/huawei_code/graph.py}

\subsection{候选方案构造与约束}
\noindent\textbf{候选方案构造与约束：}空闲核继承、规范化编号、场景化依赖合法性检查和初始分块方案。\par\nobreak
\VerbatimInput[firstline=19,lastline=96,firstnumber=19]{support/experiment/v7_fast_20260927/src/huawei_code/plan.py}
\VerbatimInput[firstline=110,lastline=150,firstnumber=110]{support/experiment/v7_fast_20260927/src/huawei_code/plan.py}

\noindent\textbf{候选变换：}局部移动、聚合、重排及共享张量的Cache候选。\par\nobreak
\VerbatimInput[firstline=13,lastline=24,firstnumber=13]{support/experiment/v7_fast_20260927/src/huawei_code/candidates.py}
\VerbatimInput[firstline=27,lastline=78,firstnumber=27]{support/experiment/v7_fast_20260927/src/huawei_code/candidates.py}
\VerbatimInput[firstline=81,lastline=97,firstnumber=81]{support/experiment/v7_fast_20260927/src/huawei_code/candidates.py}
\VerbatimInput[firstline=100,lastline=140,firstnumber=100]{support/experiment/v7_fast_20260927/src/huawei_code/candidates.py}

\subsection{轻量排序与官方评价}
\noindent\textbf{轻量排序：}互斥结构搬运、容量压力、逻辑COPY查询代理及候选排序。\par\nobreak
\VerbatimInput[firstline=31,lastline=80,firstnumber=31]{support/experiment/v7_fast_20260927/src/huawei_code/scoring.py}
\VerbatimInput[firstline=83,lastline=119,firstnumber=83]{support/experiment/v7_fast_20260927/src/huawei_code/scoring.py}
\VerbatimInput[firstline=127,lastline=206,firstnumber=127]{support/experiment/v7_fast_20260927/src/huawei_code/scoring.py}

\noindent\textbf{评价接口：}调用题面评价程序并分别返回A、B、C三问的工期、搬运和命中字段。\par\nobreak
\VerbatimInput[firstline=11,lastline=16,firstnumber=11]{support/experiment/v7_fast_20260927/src/huawei_code/official.py}
\VerbatimInput[firstline=44,lastline=58,firstnumber=44]{support/experiment/v7_fast_20260927/src/huawei_code/official.py}
\endgroup
